\documentclass[10pt,a4paper]{amsart}
\usepackage[margin=19mm]{geometry}
\usepackage{amsmath,amssymb,mathtools}
\usepackage[scr=rsfs,cal=euler]{mathalfa}
\usepackage{xfrac}
\usepackage[unicode,hidelinks,bookmarks=false]{hyperref}
\newcommand{\ie}{i.e.\ }
\newcommand{\cf}{cf.\ }
\newcommand{\resp}{respectively\ }
\newcommand{\Subsection}{\subsection}
\providecommand{\nobreakdash}{\nobreak\hbox{-}}
\setlength{\emergencystretch}{2em}
\multlinegap0pt
\usepackage{dsfont}
\usepackage{amssymb}
\newtheorem{theorem}{Theorem}
\title{Finite-time stabilization via impulse control of degenerate singular parabolic equations}
\author{Walid Zouhair, Ghita El Guermai, Ilham Ouelddris}
\date{Source: arXiv:2604.01311v1 (CC BY 4.0)}
\begin{document}
\maketitle

\noindent\emph{Source and licence.} Finite-time stabilization via impulse control of degenerate singular parabolic equations. Version arXiv:2604.01311v1 (CC BY 4.0), distributed by arXiv under the Creative Commons Attribution 4.0 International licence, http://creativecommons.org/licenses/by/4.0/, as recorded in the arXiv licence metadata for this exact version and shown on the abstract page https://arxiv.org/abs/2604.01311v1. Full licence notice: \texttt{licenses/arxiv-2604.01311-CC-BY-4.0.txt}.\par\smallskip

\noindent\textit{Editorial scope.} Counted unit: the article's theorem normop (source0.tex lines 753--768). The article's complete original proof of it is delivered with it (source0.tex lines 769--917, one \texttt{proof} environment of 149 lines). No mathematics is written, repaired or re-derived: the statement and the proof are the article's own lines, in the article's own notation, and no retained line is substituted or re-flowed.  Retained uncounted as the source-local context the proof needs: lines 82--92; lines 531--533; lines 594--600.  Nothing was omitted from any retained span, so the package carries no omission record.  No retained line contains a citation, so the wrapper carries no bibliography and no bibliography entry.  No retained line contains a figure environment, an image-inclusion command or drawing code, so no image is delivered and the package declares no asset dependency.  Editorial formatting only, in addition: an AMS \texttt{amsart} preamble that restates the article's own declarations and macros used by the retained lines, namely the environments \texttt{theorem} on separate counters, together with the standard packages those lines need.  The retained environments print in this document as Theorem 1 in order of appearance, whereas the article prints the same items with the numbering of its own full document.  The complete native source of the article as distributed in the arXiv e-print for this exact version is bundled unchanged as \texttt{sources/197655/source0.tex}.
\begin{equation}\label{1.1}
\begin{cases}\partial_t y-\left(x^\alpha y_x\right)_x-\dfrac{\mu}{x^\beta} y=0, & \text { in }(0,1) \times(0, T) \setminus\displaystyle \bigcup_{k\geq 0 }\{\tau_{k}\}, \\ 
y(\cdot, \tau_{k})=y\left(\cdot, \tau_{k}^{-}\right)+\mathds{1}_{\omega} \mathcal{L}_{k}(y(\cdot,t_{k})), &\text { in } (0,1),\\
y(1, t)=0,  & \text { on }(0, T)\\
\left\{\begin{array}{ll}
y(0, t)=0, & (W D), \\
\left(x^\alpha y_x\right)(0, t)=0, & (S D),
\end{array} \right. &\text { on }(0, T),\\
y(x, 0)=y_0(x), & \text { on }(0,1),
\end{cases}
\end{equation}

\begin{equation}\label{ImpCont}
\left\|\mathcal{L}_{k}\left(y\left(t_{k}\right)\right)\right\|_{\omega}^{2}\leq C_{2}\mathrm{e}^{2k- \frac{1}{8}\eta b^{k}} \left\|y\left(t_{0}\right)\right\|^{2},
\end{equation}

\begin{equation}\label{ACP1}
\begin{cases}
\hspace{-0.1cm} \partial_t y(t)-A y(t)=0, \qquad \qquad  \qquad(0,T)\setminus \displaystyle \bigcup _{k\geq 0 }\{\tau_{k}\}, \\
\hspace{-0.1cm} y(\tau_{k}) =y(\tau_{k}^-)+ \mathds{1}_{\omega} \mathcal{L}_{k}(y(t_k)),\\
\hspace{-0.1cm} y(0)=y_0. 
\end{cases}
\end{equation}

\begin{theorem}
The following properties are true:
    \begin{itemize}
        \item [i)] The problem $(\mathcal{P})$ has a unique minimal norm control.
        \item [ii)] The minimal norm control $(\mathcal{L}_{k}^*)_{k\geq 0}$ satisfies that $\mathcal{L}_{k}^*=0$, $k\geq 0$ if and only if the solution of \eqref{ACP1} with $\mathcal{L}_{k}(y(t_k))=0$, $k\geq 0$ satisfies \[y(T)=0.\]
        \item [iii)] Let $v_{\varepsilon,k}$ be the unique minimum of \(\mathcal{J}_{\varepsilon,k}\), then the control given by
        \begin{equation}\label{normop}
            (u_{\varepsilon,k})_{k \geq 0}= \left( \mathds{1}_{\omega}^{*} e^{(t_{k+1}-\tau_k)A} v_{\varepsilon,k}\right)_{k \geq 0},
        \end{equation}
        is a null approximate control of \eqref{1.1}, i.e
        \[\|y(T,y_0,(u_{\varepsilon,k})_{k \geq 0}) \| \leq \varepsilon\|y_0\|.\]
        Moreover, there is $(u_k)_{k \geq 0} \in \ell^2(L^2(\omega))$ such that 
        \[ u_{\epsilon,k} \to u_k \qquad \text{weakly in} \;\; \ell^2(L^2(\omega)),\]
        and $u_k$ is a norm optimal null control. That is, $u_k$ is a solution of $(\mathcal{P})$.
    \end{itemize}
\end{theorem}

\begin{proof}
    \begin{enumerate}
    \item[i)] Since the system \eqref{1.1} is null impulse controllable and the control impulse satisfies \eqref{ImpCont} , the following set is non-empty:
        \begin{equation*}
            \mathcal{F}_{ad}:=\lbrace (\mathcal{L}_{k}(y(t_k))_{k\geq 0}\in l^2(L^2(\omega)): \;\;y(T)=0 \rbrace.
        \end{equation*}
That ensures  $\mathcal{F}_{ad} \neq \varnothing$.

    On the other hand, we have $\mathcal{F}_{ad}$ is weakly closed in $l^2(L^2(\omega))$,then the problem $(\mathcal{P})$ has a minimal norm control. Now, if we consider  $(\mathcal{R}_k)_{k\geq 0}$ and $(\mathcal{S}_k)_{k\geq 0}$ as two minimal norm controls to $(\mathcal{P})$. Then we can get
    \begin{equation}\label{ega1}
    0\leq \sum\limits_{k\geq 0}\| \mathcal{R}_{k} \|_{L^2(\omega)}^2=\sum\limits_{k\geq 0}\| \mathcal{S}_{k} \|_{L^2(\omega)}^2=N^2<\infty. 
    \end{equation} 
    Furthermore, it can be verified that $\left((\mathcal{R}_k+\mathcal{S}_k)/2\right)_{k\geq 0}$ also constitutes a minimal norm control for $(\mathcal{P})$. Hence, by applying the Parallelogram low one can derive
    \begin{align}\label{ega2}
    \begin{split}
    &\sum\limits_{k\geq 0}\| (\mathcal{R}_k-\mathcal{S}_k)/2 \|_{L^2(\omega)}^2\\
    = &\dfrac{1}{2}\left( \sum\limits_{k\geq 0}\| \mathcal{R}_{k} \|_{L^2(\omega)}^2+\sum\limits_{k\geq 0}\| \mathcal{S}_{k} \|_{L^2(\omega)}^2\right) - \sum\limits_{k\geq 0}\| (\mathcal{R}_k+\mathcal{S}_k)/2 \|_{L^2(\omega)}^2\\
    = &\dfrac{1}{2}\left( \sum\limits_{k\geq 0}\| \mathcal{R}_{k} \|_{L^2(\omega)}^2+\sum\limits_{k\geq 0}\| \mathcal{S}_{k} \|_{L^2(\omega)}^2\right) - N^2.
     \end{split}
     \end{align} 
    By \eqref{ega1} and \eqref{ega2}, we obtain 
    \[ \mathcal{R}_k=\mathcal{S}_k, \quad \forall k\in \mathbb{N}. \]
    \item[ii)] This follows directly from the definition of the problem \((\mathcal{P})\).
    \item[iii)] Let us fix $k\geq 0$ and consider $v_{\varepsilon,k}$ as the unique minimum of $\mathcal{J}_{\varepsilon,k}$, i.e. 
    \[\mathcal{J}_{\varepsilon,k}(v_{\varepsilon,k}) \leq \mathcal{J}_{\varepsilon,k}(v_{\varepsilon,k}+ \lambda v), \;\;\text{for all}\;\;\lambda\in \mathbb{R}\;\;\text{and}\;\;v\in L^2(0,1).\]
    This implies that
    \begin{align*}
        &\dfrac{1}{2} \| \mathds{1}_{\omega}^{*} e^{(t_{k+1}-\tau_k)A} v_{\varepsilon,k}\|^2_{\omega} + \left\langle y_0, e^{(t_{k+1}-t_k)A} v_{\varepsilon,k} \right\rangle + \varepsilon \|y_0\|\| v_{\varepsilon,k}\|\\
        \leq &\dfrac{1}{2} \| \mathds{1}_{\omega}^{*} e^{(t_{k+1}-\tau_k)A} (v_{\varepsilon,k} + \lambda v) \|^2_{\omega} + \left\langle y_0, e^{(t_{k+1}-t_k)A} v_{\varepsilon,k} \right\rangle \\
        &+ \lambda \left\langle y_0, e^{(t_{k+1}-t_k)A} v \right\rangle + \varepsilon \|y_0\|\| v_{\varepsilon,k} + \lambda v\|\\
        \leq &\dfrac{1}{2} \| \mathds{1}_{\omega}^{*} e^{(t_{k+1}-\tau_k)A} (v_{\varepsilon,k} + \lambda v) \|^2_{\omega} + \left\langle y_0, e^{(t_{k+1}-t_k)A} v_{\varepsilon,k} \right\rangle \\
        &+ \lambda \left\langle y_0, e^{(t_{k+1}-t_k)A} v \right\rangle + \varepsilon \|y_0\|\| v_{\varepsilon,k}\| + \varepsilon \left\lvert \lambda \right\rvert \|y_0\| \|v\|,
    \end{align*}
    which enables us to write
    \begin{equation}\label{eqq12}
        \begin{aligned}
        0\leq &\dfrac{1}{2} \left(\| \mathds{1}_{\omega}^{*} e^{(t_{k+1}-\tau_k)A} (v_{\varepsilon,k}+\lambda v)\|_{\omega} \mathds{1}_{\omega}^{*} e^{(t_{k+1}-\tau_k)A} v_{\varepsilon,k}\|_{\omega} \right)\nonumber\\
        &\times \left(\| \mathds{1}_{\omega}^{*} e^{(t_{k+1}-\tau_k)A} (v_{\varepsilon,k}+\lambda v)\|_{\omega} - \| \mathds{1}_{\omega}^{*} e^{(t_{k+1}-\tau_k)A} v_{\varepsilon,k}\|_{\omega}  \right)\nonumber\\
        &+ \lambda \left\langle y_0, e^{(t_{k+1}-t_k)A} v \right\rangle +  \varepsilon \left\lvert\lambda \right\rvert \|v\| \|y_0\|.
    \end{aligned}
    \end{equation}
    Then, for every $\lambda \neq 0$ the estimate \eqref{eqq12} leads us to
    \begin{align*}
        0\leq &\dfrac{1}{2} \left(\| \mathds{1}_{\omega}^{*} e^{(t_{k+1}-\tau_k)A} (v_{\varepsilon,k}+\lambda v)\|_{\omega} + \| \mathds{1}_{\omega}^{*} e^{(t_{k+1}-\tau_k)A} v_{\varepsilon,k}\|_{\omega} \right)\nonumber\\
        &\times \dfrac{1}{\left\lvert \lambda \right\rvert} \left(\| \mathds{1}_{\omega}^{*} e^{(t_{k+1}-\tau_k)A} (v_{\epsilon,k}+\lambda v)\|_{\omega} -\| \mathds{1}_{\omega}^{*} e^{(t_{k+1}-\tau_k)A} v_{\varepsilon,k}\|_{\omega} \right)\nonumber\\
        &+ \dfrac{\lambda}{\left\lvert\lambda \right\rvert} \left\langle y_0, e^{(t_{k+1}-t_k)A} v \right\rangle +  \varepsilon \|v\| \|y_0\|.
    \end{align*}
    Passing to the limit $\lambda \to 0^{+}$ and $\lambda \to 0^{-}$ in the above inequality, one can obtain 
    \begin{equation*}
        \left\lvert \left\langle \mathds{1}_{\omega}^{*} e^{(t_{k+1}-\tau_k)A} v_{\varepsilon,k} , \mathds{1}_{\omega}^{*} e^{(t_{k+1}-\tau_k)A} v\right\rangle+ \left\langle e^{(t_{k+1}-t_k)A} y_0,  v \right\rangle \right\rvert \leq \varepsilon \| v\| \|y_0\|.
    \end{equation*}
    Then from the above estimate and \eqref{normop}, it follows that
    \begin{equation}
        \left\lvert \left\langle e^{(t_{k+1}-t_k)A} y_0 + \mathds{1}_{\omega} e^{(t_{k+1}-\tau_k)A} u_{\varepsilon,k} , v \right\rangle \right\rvert \leq \varepsilon \| v \| \|y_0\|.
    \end{equation}
    On the other hand, the solution of the impulsive system \eqref{ACP1}corresponding to the initial data $y_0$ and the control function $u_{\epsilon,k}$ over the interval $(t_k, t_{k+1})$, can be expressed as follows
    \[y(t_k,t_{k+1},y_0,u_{\varepsilon,k})= e^{(t_{k+1}-t_k)A} y_0 + \mathds{1}_{\omega} e^{(t_{k+1}-\tau_k)A} u_{\varepsilon,k},\]
    this gives that
    \[\left\lvert \left\langle y(t_k,t_{k+1},y_0,u_{\varepsilon,k}) , v \right\rangle \right\rvert \leq \varepsilon \| v \| \|y_0\|, \;\;\forall v \in L^2(0,1),\]
    which implies 
    \begin{equation}\label{eqq8}
        \| y(t_k,t_{k+1}, y_0, u_{\varepsilon,k}) \| \leq \varepsilon \|y_0\|.
    \end{equation}
    This indicates that the impulsive system \eqref{ACP1} is null approximately controllable on the interval $(t_k, t_{k+1})$. Moreover, since the system \eqref{ACP1} is finite-time stabilizable, the control satisfies
    \[ \| \mathcal{L}_k(y(t_k)) \|^2_{\omega} \leq C_2 e^{-2k} \|y_0\|^2.\]
    Now, let us prove that for any $v \in L^2(0,1)$, one has
    \begin{equation}\label{eqq13}
        \| e^{(t_{k+1}-t_k)A} v \| \leq C e^{-k} \|y_0\|\| \mathds{1}_{\omega}^{*} e^{(t_{k+1}-\tau_k)A} v \|_{\omega} + \varepsilon\|y_0\| \| v\|.
    \end{equation}
    To this end, let us consider $v \in L^2(0,1)$ and multiply \(\eqref{ACP1}_{(1)}\) by $e^{(t_{k+1}-t)A} v$ and integrate over $(0,1)$ in order to derive
    \[\left\langle \partial_t y(t) , e^{(t_{k+1}-t)A} v \right\rangle - \left\langle A \,y(t) , e^{(t_{k+1}-t)A} v \right\rangle=0.\]
    The self-adjoint nature of the operator $A$ enables us to conclude that 
    \begin{equation}\label{eqq4}
        \partial_t \left( \left\langle y (t), e^{(t_{k+1}-t)A} v \right\rangle \right) =0.
    \end{equation}
    Integrating \eqref{eqq4} over the interval $\left[t_k , \tau_k\right)$ yields
    \begin{equation}\label{eqq5}
        \left\langle y(\tau_k^{-}) , e^{(t_{k+1}-\tau_k)A} v \right\rangle - \left\langle y(t_k) , e^{(t_{k+1}-t_k)A} v \right\rangle =0.
    \end{equation}
     Similarly, integrating over the interval  $\left(\tau_k, t_{k+1}\right)$ gives
    \begin{equation}\label{eqq6}
        \left\langle y(t_{k+1}) ,  v \right\rangle - \left\langle y(\tau_k) , e^{(t_{k+1}-\tau_k)A} v \right\rangle   =0.
    \end{equation}
    Combining the results from \eqref{eqq5} and \eqref{eqq6}, and using the fact that
    \[y(\tau_k) = y(\tau^{-}_{k}) + \mathds{1}_{\omega} \mathcal{L}_k (y(t_k)),\] 
    we can write
    \begin{equation}
        \left\langle y(t_{k+1}) , v \right\rangle = \left\langle y(t_k) , e^{(t_{k+1}-t_k)A} v \right\rangle + \left\langle \mathds{1}_{\omega} \mathcal{L}_k(y(t_k)) , e^{(t_{k+1}-\tau_k)A} v \right\rangle.
    \end{equation}
    Thus, for each $k\geq 0$,
    \begin{align*}
        \| e^{(t_{k+1}-t_k)A} v \|
        &= \sup_{\| y(t_k)\| \leq 1} \left\langle e^{(t_{k+1}-t_k)A} v , y(t_k) \right\rangle\nonumber\\
        &=\sup_{\| y(t_k)\| \leq 1} \left( \left\langle y(t_{k+1}) , v \right\rangle - \left\langle \mathcal{L}_k (y(t_k)) , e^{(t_{k+1}-\tau_k ) A} v \right\rangle_{\omega} \right)\nonumber\\
        &\leq \sup_{\| y(t_k)\| \leq 1} \left( \| y(t_{k+1})\| \| v\| + \| \mathcal{L}_k (y(t_k)) \|_{\omega} \| e^{(t_{k+1}-\tau_k)A} v\|_{\omega} \right).
    \end{align*}
    Given that $\| \mathcal{L}_k (y(t_k)) \|_{\omega} \leq C e^{-k} \| y_0\|$, and using this together with \eqref{eqq8}, we can obtain the desired estimate \eqref{eqq13}. Now, let us turn back to the fact that the functional $\mathcal{J}_{\varepsilon,k}$ admits a unique minimum $v_{\epsilon,k}$, this enables us to derive
    \[\mathcal{J}^{'}_{\varepsilon,k} (v_{\varepsilon,k}) \cdot w =0, \;\;\forall w \in L^2(0,1),\]
    which implies that
    \[\left\langle \mathds{1}^{*}_{\omega} e^{(t_{k+1}-\tau_k)A} v_{\varepsilon,k}, \mathds{1}^{*}_{\omega} e^{(t_{k+1}-\tau_k)A} w \right\rangle_{\omega} + \left\langle y_0 , e^{(t_{k+1}-t_k)A} w \right\rangle_{\omega} + \varepsilon \|y_0\|\dfrac{\left\langle v_{\varepsilon,k} , w \right\rangle}{\| v_{\varepsilon,k} \|} =0,\]
    take $w=v_{\varepsilon,k}$ in the above quantity and using the estimate \eqref{eqq13} to get
    \begin{align}\label{eqq9}
       & \| \mathds{1}^{*}_{\omega} e^{(t_{k+1}-\tau_k)A} v_{\varepsilon,k} \|_{\omega}^2\nonumber\\
        &= - \left\langle y_0 , e^{(t_{k+1}-t_k)A} v_{\varepsilon,k} \right\rangle_{\omega} - \varepsilon \|y_0\| \| v_{\varepsilon,k} \|\nonumber\\
        &\leq \| y_0 \| \| e^{(t_{k+1}-t_k)A} v_{\varepsilon,k} \| - \varepsilon \|y_0\| \| v_{\varepsilon,k} \|\nonumber\\
        &\leq \left( \varepsilon \|y_0\| \| v_{\varepsilon,k} \| + C e^{-k} \| \mathds{1}^{*}_{\omega} e^{(t_{k+1}-t_k)A} v_{\varepsilon,k} \|_{\omega} \| y_0 \| \right) \| y_0\| - \varepsilon \|y_0\|\| v_{\varepsilon,k} \|\nonumber\\
        &=C e^{-k} \| \mathds{1}^{*}_{\omega} e^{(t_{k+1}-t_k)A} v_{\varepsilon,k} \|_{\omega} \| y_0 \|^2.
    \end{align}
    That allows us to give
    \[\| \mathds{1}^{*}_{\omega} e^{(t_{k+1}-\tau_k)A} v_{\varepsilon,k} \|_{\omega} \leq C e^{-k} \| y_0 \|^2.\]
    By summing over $k \geq 0$, one has
    \begin{align*}
        \sum_{k \geq 0} \| \mathds{1}^{*}_{\omega} e^{(t_{k+1}-\tau_k)A} v_{\varepsilon,k} \|_{\omega} 
        \leq C \sum_{k \geq 0} e^{-k} \| y_0 \|^2
        =\dfrac{C}{1-e^{-1}} \| y_0\|^2.
    \end{align*}
   Then, we infer that there exists a positive constant $\mathcal{K}$ such that
   \begin{equation*}
       \| (\mathds{1}^{*}_{\omega} e^{(t_{k+1}-\tau_k)A} v_{\varepsilon,k})_{k \geq 0} \|_{\ell^2(L^2(\omega))} \leq \mathcal{K} \| y_0\|^2, 
   \end{equation*}
   from \eqref{normop}, the above inequality becomes
   \begin{equation}\label{eqq10}
       \| ( u_{\varepsilon,k})_{k \geq 0} \|_{\ell^2(L^2(\omega))} \leq \mathcal{K} \| y_0\|^2.
   \end{equation}
   This allows us to deduce that there is $(u_k)_{k \geq 0} \in \ell^2(L^2(\omega))$ such that 
   \[u_{\varepsilon,k} \to u_k \;\;\text{as}\;\varepsilon \to 0 \;\;\text{weakly in}\;\;\ell^2(L^2(\omega)).\]
   By the linearity of $(y_0,v) \mapsto y(T, y_0, v)$, one has
   \begin{equation*}
       y(T,y_0,(u_{\varepsilon,k})_{k \geq 0}) \to y(T,y_0,(u_k)_{k \geq 0}) \;\;\text{as}\;\;\varepsilon\to 0\;\;\text{weakly in}\;\;L^2(0,1).
   \end{equation*}
   Moreover, one has
   \begin{align*}
       \| y(T,y_0,(u_{\varepsilon,k})_{k \geq 0})\|
       &\leq e^{\frac{T}{b^{k+1}}\delta} \| y(t_{k+1},y_0,(u_{\varepsilon,k})_{k \geq 0})\| \\
       &\leq \varepsilon \, e^{\frac{T}{b^{k+1}}\delta}\|y_0\| \to 0 \qquad\text{as}\;\;\varepsilon\to 0.
   \end{align*}
   This implies that 
   \[y(T,y_0,(u_{\varepsilon,k})_{k \geq 0})=0.\]
    Hence, it follows that \( (u_{\varepsilon,k})_{k \geq 0} \) constitutes a solution to problem \( (\mathcal{P}) \). Now, we proceed to establish the uniqueness of the solution. To do so, we assume \( u_1 \) to be an arbitrary null control for \eqref{ACP1}. Consequently,
    \begin{equation*}
       \| (u_{\varepsilon,k})_{k \geq 0} \|_{\ell^2(L^2(\omega))}\leq \| u_1\|_{\ell^2(L^2(\omega))}.
    \end{equation*}
    Since $\| (u_k)_{k \geq 0}\|_{\ell^2(L^2(\omega))} \leq \lim \inf \|(u_{\epsilon,k})_{k \geq 0}\|_{\ell^2(L^2(\omega))}$, then
   \begin{equation*}
       \| (u_k)_{k \geq 0}\|_{\ell^2(L^2(\omega))} \leq \|u_1\|_{\ell^2(L^2(\omega))}.
   \end{equation*}
   This proves that $(u_k)_{k \geq 0}$ is norm optimal control of $(\mathcal{P})$.
\end{enumerate}
\end{proof}

\end{document}
